\documentclass[10pt,a4paper]{amsart}
\usepackage[margin=19mm]{geometry}
\usepackage{amsmath,amssymb,mathtools}
\usepackage[scr=rsfs,cal=euler]{mathalfa}
\usepackage{xfrac}
\usepackage[unicode,hidelinks,bookmarks=false]{hyperref}
\newcommand{\ie}{i.e.\ }
\newcommand{\cf}{cf.\ }
\newcommand{\resp}{respectively\ }
\newcommand{\Subsection}{\subsection}
\providecommand{\nobreakdash}{\nobreak\hbox{-}}
\setlength{\emergencystretch}{2em}
\multlinegap0pt
\usepackage{dsfont}
\usepackage{enumitem}
\providecommand{\texorpdfstring}[2]{#1}
\newcommand\A{\mathrm{A}}
\newcommand\one{\mathds{1}}
\newcommand\zero{\mathbf{0}}
\newcommand\vb{\mathbf{b}}
\newcommand\vd{\mathbf{d}}
\newcommand\vf{\mathbf{f}}
\newcommand\vx{\mathbf{x}}
\newcommand\vy{\mathbf{y}}
\newcommand\vz{\mathbf{z}}
\newcommand\F{\mathbb{F}}
\newcommand\C{\mathbb{C}}
\newcommand\Z{\mathbb{Z}}
\newcommand\Q{\mathbb{Q}}
\newcommand\R{\mathbb{R}}
\renewcommand\S{\mathrm{S}}
\newcommand\tr{\mathrm{tr}}
\newcommand\Aut{\mathrm{Aut}}
\newcommand\bal{\mathrm{bal}}
\renewcommand\P{\mathcal{P}}
\newcommand\lam{\lambda}
\newcommand\lcm{\mathrm{lcm}}
\def\dotcup{\,\dot\cup\,}
\newcommand{\ontopof}[2]{\genfrac{}{}{0pt}{}{#1}{#2}}
\newcommand{\comment}[1]{}
\newcommand\aug{\fboxsep=-\fboxrule\!\!\!\fbox{\strut}\!\!\!}
\renewcommand{\theenumi}{\alph{enumi}}
\newtheorem{cdrthm}{Theorem}[section]
\theoremstyle{plain}
\newtheorem{thm}[cdrthm]{Theorem}
\newtheorem{lemma}[cdrthm]{Lemma}
\newtheorem{prop}[cdrthm]{Proposition}
\newtheorem{coro}[cdrthm]{Corollary}
\theoremstyle{definition}
\newtheorem{defi}[cdrthm]{Definition}
\newtheorem{exam}[cdrthm]{Example}
\theoremstyle{remark}
\newtheorem{rema}[cdrthm]{Remark}
\newtheorem{remas}[cdrthm]{Remarks}
\title[Base size of the symmetric group on subsets]{Exact minimum base size for symmetric actions on r-subsets and the corresponding alternating action when n is at least two r: original Theorem 1.1}
\author{Coen del Valle and Colva M. Roney-Dougal}
\date{Source: Algebraic Combinatorics 7 (2024), 959--967; DOI 10.5802/alco.370; publisher version}
\begin{document}
\maketitle
\noindent\textit{Editorial scope.} The counted unit is original Theorem 1.1 with the author's complete two-part proof, and the Section 2 and Section 3 machinery that proof uses. The hypothesis n at least 2r, the minimality of l and the defining inequalities for $m_r(l,k)$ are unchanged, every retained passage is native author text, and printed slips are kept literal without mathematical repair. Corollaries and the worked example that use the theorem are excluded and are not used to inflate the difficulty judgement.
\section{Introduction}

A \emph{base} for a permutation group $G$ acting on a set $\Omega$ is a subset $\{\alpha_1,\alpha_2,\dots, \alpha_k\}\subseteq\Omega$ with trivial pointwise stabiliser in $G$. Bases have proved to be tremendously useful in permutation group algorithms (see e.g.~\cite{ser}); for many computations the complexity is a function of the size of the base used, so it is of interest to find a smallest possible base. The size $b(G)$ of a smallest base for $G$ is called the \emph{base size} of $G$. Blaha~\cite{blaha} shows that for any positive integer $l$, and any permutation group $G$, the problem of determining whether $G$ admits a base of size at most $l$ is NP-complete. On the other hand, there are many different estimates known for $b(G)$ --- for example we can derive elementary upper and lower bounds as follows. If $\{\alpha_i\}_{i=1}^{k}$ is a base for $G$ then any group element $g\in G$ is uniquely determined by the tuple $(\alpha_i^g)_{i=1}^{k}$ and so $|G|\leq |\Omega|^{k}$. If $k=b(G)$ then $|G_{(\alpha_1,\alpha_2,\dots,\alpha_i)} : G_{(\alpha_1,\alpha_2,\dots,\alpha_{i+1})}|\geq 2$ for all $1\leq i<k$, so $2^{b(G)}\leq |G|$, and hence $(\log |G|)/(\log |\Omega|)\leq b(G)\leq \log |G|$.

In this paper we consider the primitive faithful actions of the symmetric and alternating groups. A result of Liebeck and Shalev~\cite{lie3} states that there is some absolute constant $c$ such that for $G\in\{\S_n,\mathrm{A}_n\}$ acting primitively, either $b(G)\leq c$ or up to equivalence $G$ is acting on either
\begin{itemize}
    \item[(i)] $r$-subsets of $[n]:=\{1,2,\dots, n\}$ with $2r<n$; or
    \item[(ii)] partitions of $[kl]$ into $k$ parts of size $l$ with $kl=n$.
\end{itemize}
Such actions are called \emph{standard}, and most other primitive actions of $\S_n$ and $\mathrm{A}_n$ are \emph{non-standard}; we denote the permutation groups in case (i) by $\S_{n,r}$ and $\mathrm{A}_{n,r}$, respectively.

\noindent\textit{Definition of the threshold quantity (original introduction, uncounted context).}
In this paper we completely determine $b(\S_{n,r})$ and $b(\mathrm{A}_{n,r})$. Given $l,k,r\in\mathbb{N}$ set $m_r(l,k):=\frac{1}{k}\left(lr-\sum_{i=1}^{k-1}i\binom{l}{i}\right)$. Our main result is the following.

\section*{Counted claim: original Theorem 1.1}
\begin{thm}\label{minbase}
Let $n\geq 2r$ be fixed and let $l$ be minimal such that there exists some $k\leq l+1$ satisfying $0\leq m_r(l,k)\leq\binom{l}{k}$ and $\sum_{i=0}^{k-1}\binom{l}{i}+m_r(l,k)\geq n.$ Then $b(\S_{n,r})=b(\mathrm{A}_{n+1,r})=l$.
\end{thm}

\begin{remas}
A pair $(l,k)$ satisfying the conditions of Theorem~\ref{minbase} can be seen to exist by setting $l=n$ and $k=2$. We give a natural description of the quantity $m_r(l,k)$ at the beginning of Section~\ref{minb}. In the case of the symmetric group, a similar result to Theorem~\ref{minbase} was very recently determined independently by Mecenero and Spiga~\cite{MeSp} --- both their formula and proof take a markedly different form from ours.
\end{remas}

\section{Bases and hypergraphs}\label{mach}
In this section we translate our problem into the language of hypergraphs. Define~$\mathrm{S}_{n,\leq r}$ to be the symmetric group $\mathrm{S}_n$ acting in the natural way on the set of subsets of $[n]$ of size at most $r$. Halasi~\cite{hal} shows that to determine $b(\S_{n,r})$ one may construct a minimum base for $\mathrm{S}_{n,\leq r}$.

\begin{lemma}[\cite{hal}]\label{hallem}
Fix $n\geq 2r$. Then $b(\S_{n,r})=b(\S_{n,\leq r})$.
\end{lemma}

To construct a minimum base for $\S_{n,\leq r}$ we start with a correspondence lemma which translates our bases into hypergraphs. A \emph{hypergraph} is a pair $(V,E)$ where $V$ is a set of points called \emph{vertices}, and $E$ is a set of subsets of $V$ called \emph{hyperedges}. Given a hypergraph $H=(V,E)$ and $v\in V$, define the $H$-\emph{neighbourhood} (or \emph{neighbourhood}, when clear from context) of $v$ to be the set $$N_H(v):=\{e\in E(H) : v\in e\},$$ and the \emph{degree} of $v$ to be $|N_{H}(V)|$. 

Suppose $\mathcal{A}$ is a collection of subsets of $[n]$, each of size at most $r$. If there are two points $x,y\in [n]$ which are contained in all of the same sets in $\mathcal{A}$, then each element of~$\mathcal{A}$ is fixed by the transposition $(x\, y)$, hence $\mathcal{A}$ is not a base for $\S_{n,\leq r}$. Similarly if no two such points exist then $\mathcal{A}$ is a base. We call a hypergraph \emph{irrepeating} if all vertices have distinct neighbourhoods. Therefore, a collection $\mathcal{B}$ of distinct subsets of $[n]$ each of size at most $r$ is a base for $\S_{n,r}$ if and only if the pair $([n],\mathcal{B})$ forms an irrepeating hypergraph. We will often take the view that bases \emph{are} hypergraphs, and hence refer to the hypergraph $([n],\mathcal{B})$ simply as $\mathcal{B}$.

If a hypergraph $H$ is irrepeating and has $l$ vertices, $n$ hyperedges (including possibly the empty edge), and maximum vertex degree at most $r$ then we call $H$ an \emph{$(l,n,r)$-hypergraph}. We call two bases $\mathcal{B}_1,\mathcal{B}_2$ for $\S_{n,\leq r}$ \emph{equivalent} if there exists some $\sigma\in\S_{n,\leq r}$ with $\mathcal{B}_1^\sigma=\mathcal{B}_2$.

\begin{prop}\label{corr}
Fix positive integers $l,n$, and $r$, and let $\mathcal{L}$ be the set of isomorphism classes of $(l,n,r)$-hypergraphs and $\mathcal{S}$ be the set of all equivalence classes of bases of~$\S_{n,\leq r}$ of size $l$. Then there exists a one-to-one correspondence $\rho: \mathcal{L}\to\mathcal{S}$.
\end{prop}

The correspondence $\rho$ is via a combinatorial construction known as the \emph{dual hypergraph}. Let $H$ be an irrepeating hypergraph. The \emph{dual} of $H$, denoted $H^\perp$, is the hypergraph with vertex set identified with the hyperedges of $H$, and hyperedges identified with vertices of $H$, where the incidence relations of $H^\perp$ are the reverse of those of $H$. That is $$V(H^\perp):=\{v_f : f\in E(H)\},$$ and $$E(H^\perp):=\{e_u : u\in V(H), v_{f}\in e_u\iff f\in N_H(u)\}.$$ The proof of Proposition~\ref{corr} requires a couple of easy facts on the operation $\cdot^\perp$ given in the following lemma, which follows directly from the definition.

\begin{lemma}\label{dual}
Let $\mathcal{H}$ be the space of isomorphism classes of irrepeating hypergraphs. Then $\cdot^\perp$ is an involution in $\mathrm{Sym}(\mathcal{H})$. Moreover, $$|V(H^\perp)|=|E(H)|\quad\text{and}\quad|E(H^\perp)|=|V(H)|$$ for all $H\in\mathcal{H}$.
\end{lemma}

\begin{proof}[Proof of Proposition~\ref{corr}]
Let $\mathcal{B}$ be a base for $\S_{n,\leq r}$ of size $l$. By Lemma~\ref{dual}, since $\mathcal{B}$ is irrepeating as a hypergraph it has an irrepeating dual, $H$, say, with $l$ vertices, $n$ edges, and maximum vertex degree at most $r$. That is, $H$ is an $(l,n,r)$-hypergraph. 

On the other hand, given an $(l,n,r)$-hypergraph, $K$, we deduce from Lemma~\ref{dual} that $|V(K^\perp)|=n$, $|E(K^\perp)|=l$, and the largest edge of $K^\perp$ has size at most $r$. Moreover, Lemma~\ref{dual} establishes that $K^\perp$ is irrepeating and so after relabelling the vertices of $K^\perp$ as $1,2,\dots,n$, the edges of the resulting hypergraph are indeed a base for $\S_{n,\leq r}$ of size $l$. Therefore, by Lemma~\ref{dual} if we define $\rho$ to return the edge set of the composition of the dual operation, $\cdot^\perp$, together with any such relabelling of vertices, then $\rho:\mathcal{L}\to\mathcal{S}$ is a bijection.
\end{proof}

It now follows that to determine the base size of $\S_{n,\leq r}$ --- and hence, by Lemma~\ref{hallem},~$\S_{n,r}$ --- it suffices to determine the minimum number of vertices of an irrepeating hypergraph with $n$ edges and maximum degree at most $r$. 

A hypergraph is called \emph{$k$-uniform} if its edge set consists only of edges of size $k$. A $k$-uniform hypergraph is called \emph{nearly-regular} if its degree sequence $a_1\geq a_2\geq\cdots\geq a_l$ satisfies $a_1-a_l\leq 1$. We use a result of Behrens et al.~\cite{beh} to prove the final lemma of this section.

\begin{lemma}\label{kexist}
Let $k,l,$ and $s$ be positive integers with $s\leq\binom{l}{k}$. Then there exists a nearly-regular $k$-uniform hypergraph on $l$ vertices with $s$ edges and highest degree $\left\lceil\frac{k}{l}s\right\rceil$.
\end{lemma}

\begin{proof}
If $l$ divides $ks$ then set $d=l$, otherwise let $d$ be the unique non-negative integer such that $d\left\lceil\frac{k}{l}s\right\rceil+(l-d)\left\lfloor\frac{k}{l}s\right\rfloor=ks$. Consider the sequence $(a_1,a_2,\dots,a_l)$ where $$a_i=\begin{cases}\left\lceil\frac{k}{l}s\right\rceil&\text{ for $1\leq i\leq d$}\\\left\lfloor\frac{k}{l}s\right\rfloor&\text{ for $d+1\leq i\leq l$}.\end{cases}$$ From $\binom{l}{k}\geq s$, we deduce $\binom{l-1}{k-1}=\frac{k}{l}\binom{l}{k}\geq\frac{k}{l}s$, and hence $$\binom{l-1}{k-1}=\left\lceil\binom{l-1}{k-1}\right\rceil\geq \left\lceil\frac{k}{l}s\right\rceil =a_1.$$ Therefore, by~\cite[Theorem 2.1]{beh}, there exists a $k$-uniform hypergraph $H$ with degree sequence $(a_1,a_2,\dots, a_l)$, so $H$ is nearly-regular with $s$ edges.
\end{proof}

With duality in mind, to construct a small base for $\S_{n,\leq r}$ it suffices to build an irrepeating hypergraph with some fixed number of edges, but neighbourhoods as small as possible. A natural way to do this is to successively add a smallest possible edge (in terms of set size), whilst ensuring no two vertices end up with the same neighbourhood. This is precisely how the main result of this section works. Recall $m_r(l,k)=\frac{1}{k}\left(lr-\sum_{i=1}^{k-1}i\binom{l}{i}\right)$ --- when clear from context, we omit the subscript $r$.

\begin{prop}\label{hypexist}
Fix positive integers $l,n$, and $r$ with $n\geq 2r$. Suppose there exists some $k\leq l+1$ such that $0\leq m(l,k)\leq\binom{l}{k}$, and $\sum_{i=0}^{k-1}\binom{l}{i}+m(l,k)\geq n.$
Then there exists an $(l,n,r)$-hypergraph.
\end{prop}

\begin{proof}
We construct such a hypergraph. Let $H_1$ be the unique irrepeating hypergraph on $l$ vertices with all possible edges of size at most $k-1$. 
By Lemma~\ref{kexist}, since $\lfloor m(l,k)\rfloor\leq \binom{l}{k}$ there exists some $k$-uniform hypergraph $H_2$ on $l$ vertices with exactly $\lfloor m(l,k)\rfloor$ edges, and highest degree at most $$\left\lceil \frac{k}{l}\left\lfloor \frac{1}{k}\left(lr-\sum_{i=1}^{k-1}i\binom{l}{i}\right)\right\rfloor\right\rceil \leq \left\lceil \frac{k}{l}\left( \frac{1}{k}\left(lr-\sum_{i=1}^{k-1}i\binom{l}{i}\right)\right)\right\rceil = r-\sum_{i=1}^{k-1}\binom{l-1}{i-1}.$$ Let $H$ be the irrepeating hypergraph obtained by adding the edges of $H_2$ to $H_1$. Then~$H$ has $l$ vertices, exactly $\sum_{i=0}^{k-1}\binom{l}{i}+m(l,k)\geq n$ edges, and highest degree at most $$\sum_{i=1}^{k-1}\binom{l-1}{i-1} +\left(r-\sum_{i=1}^{k-1}\binom{l-1}{i-1}\right)=r.$$ 
By arbitrarily deleting edges of size $k$ until exactly $n$ edges remain we do not increase the degree of any vertex, hence we obtain an $(l,n,r)$-hypergraph.
\end{proof}

\section{Base size of \texorpdfstring{$\S_{n,r}$}{Snr} and \texorpdfstring{$\mathrm{A}_{n,r}$}{Anr}}\label{minb}

In this section we use the tools developed in Section~\ref{mach} to deduce Theorem~\ref{minbase}. We then illustrate how the construction works in practice with a brief example, before proving a couple of corollaries. 

We first give a description of the quantity $m_r(l,k)$ in terms of bases for $\S_{n,r}$ --- it is useful to start by stating the following lemma.

\begin{lemma}\label{doubcount}
Let $n$ and $r$ be positive integers with $n\geq 2r$, and $\mathcal{B}$ a base for $\S_{n,r}$ with~$|\mathcal{B}|=l$. Then $lr=\sum_{x\in [n]}|N_{\mathcal{B}}(x)|$.
\end{lemma}

\begin{proof}
Count pairs $(B,x)$ where $x\in B\in\mathcal{B}$ in two ways.
\end{proof}

Given a base $\mathcal{B}$ as in the statement of Lemma~\ref{doubcount} and some positive integer $k$, set $$A_1:=\{x\in [n] : |N_{\mathcal{B}}(x)|< k\} \quad\text{and}\quad A_2:=\{x\in [n] : |N_{\mathcal{B}}(x)|\geq k\}.$$ We can rewrite the sum in Lemma~\ref{doubcount} as $lr=\left(\sum_{x\in A_1}|N_{\mathcal{B}}(x)|\right)+\left(\sum_{x\in A_2}|N_{\mathcal{B}}(x)|\right)$, hence $$\sum_{x\in A_2}|N_{\mathcal{B}}(x)|=lr-\left(\sum_{x\in A_1}|N_{\mathcal{B}}(x)|\right).$$ Since bases are irrepeating, it follows that there are at most $\binom{l}{i}$ distinct neighbourhoods of size $i$, hence $lr-\left(\sum_{x\in A_1}|N_{\mathcal{B}}(x)|\right)$ is minimally $$lr-\sum_{i=1}^{k-1}i\binom{l}{i}=km_r(l,k).$$ On the other hand $k|A_2|\leq\sum_{x\in A_2}|N_{\mathcal{B}}(x)|$. Therefore $m_r(l,k)$ estimates the minimum number of points of $[n]$ which have $\mathcal{B}$-neighbourhoods of size at least $k$.


\section*{Counted proof: the author's complete proof of Theorem 1.1}
We now proceed with the proof of the symmetric group case of Theorem~\ref{minbase}.

\begin{proof}[Proof of Theorem~\ref{minbase} for $\S_n$]
Let $\mathcal{B}$ be a minimum base for $\S_{n,r}$ so that $b:=b(\S_{n,r})=|\mathcal{B}|$. Let $l$ be minimal satisfying the conditions of Proposition~\ref{hypexist}. Then there exists an $(l,n,r)$-hypergraph, $H$, say, and by Proposition~\ref{corr}, $\rho(H)$ is a base of size $l$ for~$\S_{n,\leq r}$. By the minimality of $b$ and Lemma~\ref{hallem} we deduce $l\geq b$. 

Now, let $h\leq b+1$ be maximal such that $0\leq\left(br-\sum_{i=1}^{h-1}i\binom{b}{i}\right)$ (note $h$ exists since the inequality holds with $h=1<b+1$). It follows from the maximality of $h$ that $$0\leq\frac{1}{h}\left(br-\sum_{i=1}^{h-1}i\binom{b}{i}\right)\leq\binom{b}{h}.$$ From Lemma~\ref{doubcount} and the subsequent discussion we deduce that $$br\geq\left(\sum_{i=0}^{h-1}i\binom{b}{i}\right)+h\left(n-\sum_{i=0}^{h-1}\binom{b}{i}\right),$$ where the second summand describes the fact that any point not contributing to the first summand has neighbourhood of size at least $h$. Rearranging gives $$\frac{1}{h}\left(br-\sum_{i=0}^{h-1}i\binom{b}{i}\right)\geq n-\sum_{i=0}^{h-1}\binom{b}{i}$$ and so $\sum_{i=0}^{h-1}\binom{b}{i}+m(b,h)\geq n$. Thus all conditions of Proposition~\ref{hypexist} are satisfied. By definition $l$ is the smallest positive integer satisfying the conditions of Proposition~\ref{hypexist} and so $l\leq b$, hence $l=b$ as desired.
\end{proof}

We now consider the alternating case.

\begin{proof}[Proof of Theorem~\ref{minbase} for $\mathrm{A}_n$]
We prove the result by showing that $b(\mathrm{A}_{n+1,r})$ is equal to $b(\mathrm{S}_{n,r}).$ Let $\mathcal{B}$ be a base for $\S_{n,r}$ of size $b(\S_{n,r})$. Consider $\mathcal{B}$ as a collection of $r$-subsets of $[n+1]$, with $n+1$ having empty neighbourhood. Since $\mathcal{B}$ is a base for $\S_{n,r}$, at most one element of $[n]$ has an empty $\mathcal{B}$-neighbourhood, hence at most two elements of $[n+1]$ do, with all others distinct. But being a base for $\mathrm{A}_{n+1,r}$ is equivalent to having at most two points with equal neighbourhoods, hence $\mathcal{B}$ is a base for $\mathrm{A}_{n+1,r}$ of size $b(\S_{n,r})$. This shows that $b(\mathrm{A}_{n+1,r})\leq b(\S_{n,r})$.

On the other hand, let $\mathcal{C}$ be a base for $\mathrm{A}_{n+1,r}$. If there are two points $x$ and $y$ with the same $\mathcal{C}$-neighbourhood, then assume without loss of generality that $x$ is $n+1$. Let $\mathcal{C}'$ be obtained from $\mathcal{C}$ by deleting $n+1$ from all $r$-sets in $\mathcal{C}$. Then $\mathcal{C}'$ is a collection of subsets of $[n]$ of size at most $r$ such that each element of $[n]$ has a distinct neighbourhood, that is, a base for $\S_{n,\leq r}$. Thus $b(\mathrm{A}_{n+1,r})\geq b(\S_{n,\leq r})$, and the result follows from Lemma~\ref{hallem}.
\end{proof}

\noindent\textit{Excluded material (recorded, not claimed).} Corollary 1.3, Example 3.2 with its figure, Corollaries 3.3 and 3.4 with their proofs, Remark 3.5 and the closing remark are not part of this unit and are not reproduced; they either restate known results from the literature or apply the counted theorem to special ranges, and the original figure is an external illustration that is not redistributed. The historical paragraphs of the introduction are likewise omitted, while the two external results actually used by the counted proof, Lemma 2.1 of Halasi and the degree-sequence theorem of Behrens et al., are quoted with their exact published citations.
\begin{thebibliography}{99}
\bibitem[1]{beh} Sarah Behrens, Catherine Erbes, Michael Ferrara, Stephen G. Hartke, Benjamin Reiniger, Hannah Spinoza, and Charles Tomlinson, ``New results on degree sequences of uniform hypergraphs'', \emph{Electron. J. Combin.} \textbf{20} (2013), no. 4, article no. 14 (18 pages), https://doi.org/10.37236/3414.
\bibitem[2]{blaha} Kenneth D. Blaha, ``Minimum bases for permutation groups: the greedy approximation'', \emph{J. Algorithms} \textbf{13} (1992), no. 2, 297--306, https://doi.org/10.1016/0196-6774(92)90020-D.
\bibitem[8]{hal} Zolt\'an Halasi, ``On the base size for the symmetric group acting on subsets'', \emph{Studia Sci. Math. Hungar.} \textbf{49} (2012), no. 4, 492--500, https://doi.org/10.1556/SScMath.49.2012.4.1222.
\bibitem[9]{lie3} Martin W. Liebeck and Aner Shalev, ``Simple groups, permutation groups, and probability'', \emph{J. Amer. Math. Soc.} \textbf{12} (1999), no. 2, 497--520, https://doi.org/10.1090/S0894-0347-99-00288-X.
\bibitem[10]{MeSp} Giovanni Mecenero and Pablo Spiga, ``A formula for the base size of the symmetric group in its action on subsets'', \emph{Australas. J. Combin.} \textbf{88} (2024), no. 2, 244--255.
\bibitem[12]{ser} \'Akos Seress, \emph{Permutation group algorithms}, Cambridge Tracts in Mathematics, vol. 152, Cambridge University Press, Cambridge, 2003, https://doi.org/10.1017/CBO9780511546549.
\end{thebibliography}
\end{document}
